\documentclass[11pt]{amsart}

\usepackage{amsmath}
\usepackage{amssymb}
\usepackage{amsthm}
%\usepackage{psfig}

\begin{document}
\baselineskip=15pt
\textheight=8.4in
\parindent=0pt
\def\sk {\hskip .5cm}
\def\skv {\vskip .12cm}
\def\cos {\mbox{cos}}
\def\sin {\mbox{sin}}
\def\tan {\mbox{tan}}
\def\intl{\int\limits}
\def\lm{\lim\limits}
\newcommand{\frc}{\displaystyle\frac}
\def\xbf{{\mathbf x}}
\def\fbf{{\mathbf f}}
\def\gbf{{\mathbf g}}

\def\Ker{{\rm Ker\,}}
\def\phi{\varphi}

\def\dbA{{\mathbb A}}
\def\dbB{{\mathbb B}}
\def\dbC{{\mathbb C}}
\def\dbD{{\mathbb D}}
\def\dbE{{\mathbb E}}
\def\dbF{{\mathbb F}}
\def\dbG{{\mathbb G}}
\def\dbH{{\mathbb H}}
\def\dbI{{\mathbb I}}
\def\dbJ{{\mathbb J}}
\def\dbK{{\mathbb K}}
\def\dbL{{\mathbb L}}
\def\dbM{{\mathbb M}}
\def\dbN{{\mathbb N}}
\def\dbO{{\mathbb O}}
\def\dbP{{\mathbb P}}
\def\dbQ{{\mathbb Q}}
\def\dbR{{\mathbb R}}
\def\dbS{{\mathbb S}}
\def\dbT{{\mathbb T}}
\def\dbU{{\mathbb U}}
\def\dbV{{\mathbb V}}
\def\dbW{{\mathbb W}}
\def\dbX{{\mathbb X}}
\def\dbY{{\mathbb Y}}
\def\dbZ{{\mathbb Z}}


\def\char{{\rm char}}
\def\wt{{\rm wt}}
\def\Aut{{\rm Aut}}
\def\Gal{{\rm Gal}}
\def\Hom{{\rm Hom}}
\def\End{{\rm End}}
\def\Mat{{\rm Mat}}
\def\Irr{\mathrm {Irr}}
\def\Im{{\rm Im}}
\def\GL{{\rm GL}}
\def\SL{{\rm SL}}
\def\SNF{{\rm SNF}}
\def\tr {{\rm tr}}
\def\rk{{\rm rk}}
\def\lam{{\lambda}}

\def\la{{\langle}}
\def\ra{{\rangle}}

\bf
\begin{center}
{Homework \# 10, due on Canvas by 11:59pm on Fri, Apr 17th.}\rm
\end{center}
\vskip .2cm
\rm

{\bf Plan for the next 2 classes:} Galois correspondence (14.2 in DF, Lectures 20-21 from Spring 2010 and Lectures 16-17 from Spring 2021). Galois theory of finite fields (14.3 in DF, Lecture 22 from Spring 2010 and Lectures 14 and 18 from Spring 2021). Cyclotomic extensions, time permitting (14.5 in DF and Lecture 18 from Spring 2021).
\vskip .1cm

%{\bf Note on hints:} Some hints are given at the end of the assignment, each on a separate page.
%Problems (or parts of problems) for which a hint is available at the end are marked with *.


\skv
{\bf Problem 1:} \rm
Prove the interesting part of Corollary 18.7 form Spring 2010 notes: if $K/L/F$ is a tower of algebraic extensions and $K/L$ and $L/F$ are separable, then $K/F$ is separable (see the notes for a hint). Give a detailed argument.
\skv

\sk Before solving Problem 2 read the statement and proof of the Primitive Element Theorem (either Lecture 18 in Spring 2010 notes or Lecture 13 in Spring 21 notes).

\skv
{\bf Problem 2:} \rm Let $p_1,\ldots, p_n$ be distinct primes. Let $\alpha_i=\sqrt[p_i]{p_i}$ and $\alpha=\sum\limits_{i=1}^n \alpha_i$. Prove that $\dbQ(\alpha_1,\ldots,\alpha_n)=\dbQ(\alpha)$. {\bf Hint:} The general approach should be similar to the examples discussed in Lecture 18 from 2010 and Lecture 13 from Spring 21; however, the proof will involve additional technicalities.


\skv
{\bf Problem 3:} \rm Let $p$ be a prime, $K=\mathbb F_p(s,t)$, the field of rational functions over $\dbF_p$ in two variables, and
let $F=\mathbb F_p(s^p,t^p)$. Prove that the extension $K/F$ cannot be generated by a single element. {\bf Hint:} First compute
$[K:F]$.
\skv

\skv
{\bf Problem 4:} Let $f(x)\in \dbQ[x]$ be an irreducible polynomial
of degree $n$, and let $K\subset \dbC$ be the splitting field for $f(x)$ over $\dbQ$.
Label the roots of $f(x)$ by $1,\ldots, n$ (in some order),
and let $\iota: \Gal(K/\dbQ)\to S_n$ be the associated embedding.

\begin{itemize}

\item[(a)] Prove that $|\Gal(K/\dbQ)|$ is a multiple of $n$ and divides $n!$

 \item[(b)] Assume $f(x)$ has at least one non-real root. Prove
that the complex conjugation is an element of $\Gal(K/\dbQ)$
of order $2$.

\item[(c)] Assume that $f(x)$ has precisely two non-real roots. Prove
that the image of the complex conjugation under the embedding $\iota$
is a transposition.

\item[(d)] Suppose that $n=\deg(f)$ is prime and again assume that $f(x)$ has precisely
two non-real roots. Prove that $\Gal(K/\dbQ)$ is isomorphic to $S_n$.
{\bf Hint:} $\Gal(K/\dbQ)$ must contain an element of order $n$ (why?)
\end{itemize}
\newpage
Before doing problems 5-7 below read \S~13.6 in DF.

\skv
{\bf Problem 5:} \rm 
Let $p$ be a prime, $n\geq 2$ an integer, $f(x)=x^n-p$, and let
$K\subset \dbC$ be the splitting field for $f(x)$ over $\dbQ$.
\begin{itemize}
\item[(a)] Prove that $K=\dbQ(\sqrt[n]{p},\omega_n)$ where
$\omega_n=e^{2\pi i/n}$.
\item[(b)] Prove that $[K:\dbQ]\leq n\phi(n)$ where $\phi$
is the Euler function. Then describe the elements of $\Gal(K/\dbQ)$ by their actions on $\sqrt[n]{p}$
and $\omega_n$.
\item[(c)] Let $G=\dbZ/n\dbZ\rtimes (\dbZ/n\dbZ)^{\times}$ where 
$(\dbZ/n\dbZ)^{\times}$
acts on $\dbZ/n\dbZ$ by multiplication. Prove that $\Gal(K/\dbQ)$ is isomorphic
to a subgroup of $G$ and describe an explicit embedding.
\item[(d)] Let $p=3$ and $n=12$. Prove that inequality in (b)
is strict and find $[K:\dbQ]$. {\bf Hint:} Compute $\omega_{12}$
explicitly.
\end{itemize}


\skv
{\bf Problem 6:} Problem 4 in [DF, p.555].
\skv
{\bf Problem 7:} Problem 5 in [DF, p.555].
\skv

\end{document}
